\chapter{Introduction}

Data breaches have grown exponentially, with over 1800 reported breaches in the first half of 2026, putting it on course to set a new annual record. AI has accounted for an increase of over 56\% of breaches from the previous year and with over 6 billion Internet users as of April,2026, placing the need for personal data governance at an all-time high at the moment. While newer regulations are constantly being adapted to ensure compliance, that hasn't hindered malicious actors from playing their hand and thereby affecting millions worldwide. This need for personal data privacy is further emphasized when large-scale infrastructural breaches such as the 2026 Canvas data breach occur, compromising 275 million users and their Personally Identifiable Information (PII) worldwide. While laws such as the General Data Protection Act (GDPR) in the E.U, the California Consumer Privacy Act (CCPA) in the USA and the Personal Information Protection Law (PIPL) in China have reinforced user data protection, enabling stricter regulations and tighter controls worldwide, their efficiency in dealing with the blurry lines has often come under scrutiny, especially with respect to their usage, adoption and understanding by lay-people. 

With a population of over 1.47 billion people, online penetration of 70\% as of October, 2025 and a new Digital Personal Data Protection Act (DPDPA) in partial enforcement from November, 2025, India is heading towards being a landmark example of digital privacy enforcement in de-centralized environments and also a certain opportunity for large-scale breaches with an existing evidenced past. This new act which was first passed in 2023, aims to reinforce India's digital landscape, focusing on the user's autonomy over their digital personal data, while providing lawful practices for organizations to process such data. Prior works on the Act discuss about the user's understanding of its laws and the need for such regulations. However, past research in aforementioned established contexts has shown that when such regulations are implemented without the consideration of the general population's needs and understanding, it could lead to frustration over certain practices, loopholes for organizations to bypass and overall skepticism on the implementation. 

Considering that the DPDPA is heading towards a complete enforcement and compliance from all the organizations from November, 2026, it is therefore important that the users' perspectives on their personal data privacy, sensitivities of various personal data items and their processing be understood to get a holistic perspective of the actual needs that need to be fulfilled and protected by the Act. This is especially important given that research has found users to emphasize the importance of their financial data in the past but haven't paid much heed to any other forms of personal data, which could potentially lead to their profiling and misuse by exploitative companies or practices. Additionally, given the recency of the Act, and its provisions to provide the users with the option to request collected data and consent revocal for its processing, it is especially important to understand their perspectives on such personal data processing and their perceived control throughout the process. While many legal experts have discussed and dissected the Act's vague definitions and their implications, an understanding on the possible gap between the user's perceptions of such definitions and process to the actual provisions of the Act would therefore establish a baseline comparative differentiation to rely upon. 

To address these gaps, an exploratory survey of 239 primarily educated participants with potential familiarity of online spaces was conducted. The analysis focused on the end user's perspectives of their general privacy perceptions of their personal data, their perceived sensitivity of the various personal data items and their perceived control over such data's processing. The study focuses on providing an initial understanding of the user's perceptions over this entire process, thereby understanding their (mis)alignment with the provisions of the Act. To do so, the following research questions were addressed:

\begin{enumerate}
    \item What are Indian Internet users’ perceptions about the privacy-sensitive nature of different types of personal data?
    \item What are Indian Internet users' beliefs regarding their control over different types of personal data and its processing?
\end{enumerate}
